\documentclass[runningheads]{llncs}

\usepackage[T1]{fontenc}
\usepackage[utf8]{inputenc}
\usepackage[spanish,es-noshorthands,es-tabla]{babel}
\usepackage{lmodern}
\usepackage{amsmath,amssymb}
\usepackage{booktabs}
\usepackage{graphicx}
\usepackage{xcolor}
\usepackage{listings}
\usepackage{tikz}
\usetikzlibrary{positioning,arrows.meta}
\usepackage{pgfplots}
\pgfplotsset{compat=1.18}
\usepackage[hidelinks]{hyperref}
\usepackage{url}

% ---------------------------------------------------------------------------
% Enlace al repositorio: reemplazar USUARIO por el usuario de GitHub
\newcommand{\repo}{https://github.com/USUARIO/misioneros-canibales-busqueda}
% ---------------------------------------------------------------------------

\definecolor{codebg}{RGB}{248,248,248}
\definecolor{codekw}{RGB}{0,70,140}
\definecolor{codecm}{RGB}{110,110,110}
\definecolor{solucion}{RGB}{200,40,40}
\lstset{
  language=Python, basicstyle=\ttfamily\scriptsize, keywordstyle=\color{codekw}\bfseries,
  commentstyle=\color{codecm}\itshape, stringstyle=\color{teal}, backgroundcolor=\color{codebg},
  frame=single, rulecolor=\color{black!20}, numbers=left, numberstyle=\tiny\color{codecm},
  xleftmargin=2em, framexleftmargin=1.5em, showstringspaces=false, columns=fullflexible,
  breaklines=true, captionpos=b,
  literate={á}{{\'a}}1 {é}{{\'e}}1 {í}{{\'i}}1 {ó}{{\'o}}1 {ú}{{\'u}}1 {ñ}{{\~n}}1
           {Á}{{\'A}}1 {É}{{\'E}}1 {Í}{{\'I}}1 {Ó}{{\'O}}1 {Ú}{{\'U}}1
}
\renewcommand{\lstlistingname}{Listado}
\AtBeginDocument{\renewcommand{\thelstlisting}{\arabic{lstlisting}}}

\newcommand{\Sr}{S_{r}}
\newcommand{\st}[3]{(#1,#2,#3)}

\begin{document}

\title{Formulación y Búsqueda Ciega en Espacios de Estados:\\
El Problema de los Misioneros y Caníbales}
\titlerunning{Búsqueda ciega en Misioneros y Caníbales}

\author{Christopher Pulupa}
\authorrunning{C. Pulupa}
\institute{Facultad de Ingeniería de Sistemas, Escuela Politécnica Nacional,\\
Quito, Ecuador\\
Tarea Individual DG\_S4 --- Inteligencia Artificial}

\maketitle

\begin{abstract}
Se formaliza el problema de los misioneros y caníbales como un problema de
búsqueda en espacio de estados $(S, s_0, A, T, G)$ y se implementan en Python
las búsquedas no informadas en anchura (BFS, frontera FIFO) y en profundidad
(DFS, frontera LIFO), junto con A$^*$ como referencia informada. En la instancia
clásica ($N=3$, bote de capacidad $K=2$) ambos algoritmos ciegos encuentran un
camino de 11 cruces expandiendo 14 (BFS) y 11 (DFS) nodos en unas decenas de
microsegundos. El análisis empírico muestra que este resultado no contradice la
teoría sino que la confirma: la restricción de seguridad reduce el espacio a
16 estados alcanzables con un factor de ramificación efectivo
$b^*\approx 1{,}04$, y en ese grafo todos los caminos simples a la meta miden
exactamente 11. Al generalizar a $N$ y $K$ mayores, DFS deja de ser óptimo
(llega a duplicar la longitud óptima) y, al retirar la detección de estados
repetidos, la frontera de BFS crece a 13\,835 nodos frente a 5 de DFS,
reproduciendo la brecha $O(b^d)$ frente a $O(bm)$.
\keywords{Búsqueda no informada \and BFS \and DFS \and Espacio de estados \and
Factor de ramificación efectivo \and Misioneros y caníbales.}
\end{abstract}

% ===========================================================================
\section{Introducción}

El problema de los misioneros y caníbales es uno de los ejemplos canónicos de la
Inteligencia Artificial para estudiar la \emph{formulación} de problemas: Amarel
\cite{amarel1968} lo usó para mostrar que la elección de la representación
determina la dificultad de la búsqueda. Tres misioneros y tres caníbales deben
cruzar un río en un bote de dos plazas sin que, en ninguna orilla, los caníbales
superen en número a los misioneros presentes.

Siguiendo el marco de agentes de resolución de problemas de Russell y Norvig
\cite{russell2020}, el agente (i) formula el objetivo y el problema, (ii) busca
una secuencia de acciones en el modelo y (iii) la ejecuta. Este trabajo cubre las
dos primeras fases: la Sección~\ref{sec:formulacion} define el problema con
rigor matemático, la Sección~\ref{sec:implementacion} describe su abstracción a
código, la Sección~\ref{sec:experimentos} presenta el \emph{benchmarking} y la
Sección~\ref{sec:discusion} contrasta los resultados con la teoría asintótica.

% ===========================================================================
\section{Formulación teórica}\label{sec:formulacion}

Se trabaja con la generalización de $N$ misioneros, $N$ caníbales y un bote de
capacidad $K$; la instancia pedida es $N=3$, $K=2$. Un problema de búsqueda es la
tupla $\mathcal{P} = (S, s_0, A, \textsc{Aplicable}, T, G, c)$.

\subsubsection{Estados.}
Basta con describir la orilla izquierda, porque la derecha se deduce por
complemento. Un estado es la tupla
\begin{equation}
  s = (m, c, b), \qquad m, c \in \{0,\dots,N\},\; b\in\{0,1\},
\end{equation}
donde $m$ y $c$ son los misioneros y caníbales en la orilla izquierda y $b=1$
si el bote está en ella. Las cantidades en la orilla derecha son
$\bar m = N-m$ y $\bar c = N-c$. El conjunto de estados \emph{válidos} es
\begin{equation}
  S = \bigl\{ (m,c,b) \;:\; (m = 0 \lor m \ge c) \;\land\; (\bar m = 0 \lor \bar m \ge \bar c) \bigr\}.
  \label{eq:seguridad}
\end{equation}
Los pares $(m,c)$ que cumplen \eqref{eq:seguridad} son los de $m=0$, los de
$m=N$ y los de la diagonal $m=c$, es decir $3N+1$ pares; por tanto
$|S| = 2(3N+1)$, que para $N=3$ da $|S|=20$ frente a las $4\cdot 4\cdot 2=32$
combinaciones sin restricción. Cuatro de ellos,
$\st{0}{0}{1}$, $\st{3}{3}{0}$, $\st{0}{3}{0}$ y $\st{3}{0}{1}$, son
inalcanzables desde $s_0$ (el bote no puede estar en una orilla vacía de la que
nadie lo trajo), de modo que el espacio alcanzable tiene $|\Sr| = 16$ estados.

\subsubsection{Estado inicial y estado objetivo.}
\begin{equation}
  s_0 = (N, N, 1) = (3,3,1), \qquad G = \{(0,0,0)\}.
\end{equation}

\subsubsection{Acciones.}
Una acción es la carga del bote, $a=(\delta_m, \delta_c)$, con
\begin{equation}
  A = \{(\delta_m,\delta_c)\in\mathbb{N}^2 : 1 \le \delta_m+\delta_c \le K\},
  \qquad |A| = \tfrac{K(K+3)}{2}.
\end{equation}
Para $K=2$, $A=\{(2,0),(1,1),(0,2),(1,0),(0,1)\}$. La acción es
\emph{aplicable} en $s$ si las personas están en la orilla donde se encuentra el bote:
\begin{equation}
  \textsc{Aplicable}(s,a) \iff
  \begin{cases}
    \delta_m \le m \;\land\; \delta_c \le c, & b = 1,\\
    \delta_m \le N-m \;\land\; \delta_c \le N-c, & b = 0.
  \end{cases}
\end{equation}

\subsubsection{Modelo de transición.}
El bote cambia de orilla y lleva consigo la carga:
\begin{equation}
  T\bigl((m,c,b),(\delta_m,\delta_c)\bigr) =
  \begin{cases}
    (m-\delta_m,\; c-\delta_c,\; 0), & b=1 \text{ (ida)},\\
    (m+\delta_m,\; c+\delta_c,\; 1), & b=0 \text{ (vuelta)},
  \end{cases}
\end{equation}
y la función de sucesores conserva solo los resultados seguros:
\begin{equation}
  \textsc{Suc}(s) = \{(a, T(s,a)) : \textsc{Aplicable}(s,a) \land T(s,a)\in S\}.
\end{equation}
Toda transición es reversible, por lo que el grafo de estados es no dirigido.

\subsubsection{Test objetivo y costo.}
$\textsc{Objetivo}(s) \iff s = (0,0,0)$. Cada cruce cuesta $c(s,a,s')=1$, así que
el costo de un camino es su número de acciones y el camino óptimo es el más corto.

La Fig.~\ref{fig:grafo} muestra el grafo completo de $\Sr$ para $N=3$, $K=2$,
ordenado por profundidad desde $s_0$. Tiene 17 aristas y un grado medio de
$2{,}125$, muy por debajo del máximo teórico $|A|=5$: la restricción
\eqref{eq:seguridad} poda la mayor parte de las acciones.

\begin{figure}[t]
\centering
\resizebox{\textwidth}{!}{%
\begin{tikzpicture}[
  every node/.style={draw, rounded corners=2pt, font=\footnotesize\ttfamily,
                     inner sep=2pt, minimum width=7mm, minimum height=5mm, fill=white},
  ini/.style={fill=blue!12}, fin/.style={fill=green!18},
  sol/.style={draw=solucion, line width=1.3pt},
  arista/.style={draw=black!45}]
  \def\dx{0.97}
  \node[ini] (331) at (0*\dx,0) {331};
  \node (220) at (1*\dx, 0.8) {220};
  \node (310) at (1*\dx, 0)    {310};
  \node (320) at (1*\dx,-0.8) {320};
  \node (321) at (2*\dx,0) {321};
  \node (300) at (3*\dx,0) {300};
  \node (311) at (4*\dx,0) {311};
  \node (110) at (5*\dx,0) {110};
  \node (221) at (6*\dx,0) {221};
  \node (020) at (7*\dx,0) {020};
  \node (031) at (8*\dx,0) {031};
  \node (010) at (9*\dx,0) {010};
  \node (111) at (10*\dx, 0.5) {111};
  \node (021) at (10*\dx,-0.5) {021};
  \node[fin] (000) at (11*\dx,0) {000};
  \node (011) at (12*\dx,0) {011};
  \foreach \a/\b in {331/310,331/320,310/321,010/021,021/000,000/011}
    \draw[arista] (\a) -- (\b);
  \foreach \a/\b in {331/220,220/321,321/300,300/311,311/110,110/221,221/020,
                     020/031,031/010,010/111,111/000}
    \draw[sol] (\a) -- (\b);
  \foreach \d in {0,...,12}
    \node[draw=none, fill=none, font=\tiny\rmfamily, text=black!55] at (\d*\dx,-1.35) {$d{=}\d$};
\end{tikzpicture}}
\caption{Grafo del espacio de estados alcanzable ($|\Sr|=16$, 17 aristas) para
$N=3$, $K=2$. La etiqueta \texttt{mcb} denota el estado $(m,c,b)$ y la columna indica la profundidad mínima
desde $s_0$. En rojo, el camino de 11 cruces encontrado por BFS. Las únicas
bifurcaciones son la primera y la última jugada; $\st{3}{2}{0}$ y
$\st{0}{1}{1}$ son callejones sin salida.}
\label{fig:grafo}
\end{figure}

% ===========================================================================
\section{Implementación}\label{sec:implementacion}

El código está en Python~3 sin dependencias externas y se organiza como un agente
de resolución de problemas con responsabilidades separadas:
el módulo \texttt{src/problema.py} encapsula la formulación $\mathcal P$, con
los estados como tuplas inmutables \texttt{(m\_izq,~c\_izq,~bote\_izq)} que sirven
de clave en tablas hash; el módulo \texttt{src/busqueda.py} contiene algoritmos
que solo conocen la interfaz del problema (\texttt{estado\_inicial},
\texttt{sucesores} y \texttt{es\_objetivo}); y el script \texttt{benchmark.py}
ejecuta los experimentos. El repositorio incluye pruebas
unitarias que verifican la formulación, la validez de cada camino y la
consistencia de la heurística.

\subsubsection{Poda en la generación.}
La restricción \eqref{eq:seguridad} se aplica dentro de \textsc{Suc}
(Listado~\ref{lst:suc}): un estado inseguro nunca se crea ni entra a la frontera.
Es la forma más simple de propagación de restricciones y evita gastar
expansiones en estados que el test objetivo rechazaría más tarde.

\begin{lstlisting}[caption={Función de sucesores con poda de estados inseguros.},label={lst:suc}]
def sucesores(self, s):
    for a in self.acciones:                    # (dm, dc), 1 <= dm + dc <= K
        if self.es_aplicable(s, a):
            s2 = self.resultado(s, a)          # modelo de transición T(s, a)
            if self.es_valido(s2):             # restricción de seguridad
                yield a, s2
\end{lstlisting}

\subsubsection{BFS y DFS.}
Ambos son búsquedas de \emph{grafo}: guardan los estados ya vistos para no
expandir dos veces el mismo estado. BFS usa \texttt{collections.deque} como cola
FIFO y marca un estado como alcanzado al generarlo; DFS usa una lista como pila
LIFO y marca al expandir, apilando los hijos en orden inverso para respetar el
orden de las acciones (Listado~\ref{lst:bfs}). Sin el conjunto de explorados,
DFS no termina: la acción $(0,1)$ seguida de su inversa genera un ciclo
infinito. En ambos el test objetivo se aplica al extraer el nodo, con lo que
los contadores de nodos expandidos son directamente comparables.

\begin{lstlisting}[caption={Núcleo de BFS; DFS cambia \texttt{popleft()} por \texttt{pop()}.},label={lst:bfs}]
frontera, alcanzados = deque([Nodo(s0)]), {s0}
while frontera:
    nodo = frontera.popleft()                  # FIFO
    if problema.es_objetivo(nodo.estado):
        return nodo                            # el camino se reconstruye por los padres
    expandidos += 1
    for accion, s2 in problema.sucesores(nodo.estado):
        if s2 not in alcanzados:
            alcanzados.add(s2)
            frontera.append(Nodo(s2, nodo, accion, nodo.profundidad + 1))
\end{lstlisting}

\subsubsection{Variantes de árbol y A$^*$.}
Para aislar el efecto de la estrategia sobre la memoria se implementaron dos
variantes sin conjunto global de explorados: BFS de árbol, y DFS de árbol que solo
evita repetir estados del camino actual, con memoria $O(bm)$. Además se incluyó
A$^*$ \cite{hart1968} con una heurística obtenida por \emph{relajación}
\cite{pearl1984}: se ignora la restricción \eqref{eq:seguridad} y se cuenta el
mínimo de cruces para pasar $n=m+c$ personas. Cada viaje de ida y vuelta traslada
netamente $K-1$ personas y el último traslada $K$; si el bote está a la derecha, primero debe volver con al menos una persona. De ahí
\begin{equation}
  h_1(n) = \begin{cases} 0, & n=0\\ 1, & 0<n\le K\\ 2\left\lceil \frac{n-K}{K-1}\right\rceil+1, & n>K \end{cases}
\end{equation}
\begin{equation}
  h(m,c,b) = \begin{cases} h_1(n), & b=1\\ 1 + h_1(n+1), & b=0,\ n>0\\ 0, & b=0,\ n=0 \end{cases}
\end{equation}
Al ser el costo exacto de un problema relajado, $h$ es admisible y consistente;
las pruebas lo comprueban sobre todas las aristas ($h(s)\le 1+h(s')$, $h(s_G)=0$).
Para $s_0$ con $K=2$ se obtiene $h(s_0)=9$, frente al costo real $h^*(s_0)=11$.

% ===========================================================================
\section{Benchmarking}\label{sec:experimentos}

\subsubsection{Protocolo.}
Las pruebas se ejecutaron con Python~3.13.18 en Linux x86\_64. El tiempo
reportado es la \emph{mediana} de 2\,000 ejecuciones (10 para BFS de árbol)
medidas con el reloj de alta resolución \texttt{perf\_counter} de Python, lo que amortigua el ruido
del sistema operativo en tiempos de microsegundos. La memoria se reporta de dos
formas: en nodos, como el máximo de $|F|+|E|$ (frontera más explorados), y en
bytes, como el pico medido con \texttt{tracemalloc} en una ejecución separada
para no sesgar el tiempo. El factor de ramificación efectivo $b^*$ se obtiene
resolviendo por bisección $N_{\text{gen}}+1 = \sum_{i=0}^{d} (b^*)^i$, donde
$N_{\text{gen}}$ son los nodos generados y $d$ la profundidad de la solución
\cite{russell2020}. Todo se reproduce con \texttt{python3 benchmark.py}.

\subsubsection{Resultados de la instancia clásica.}
La Tabla~\ref{tab:principal} es la comparación que pide el enunciado y la
Tabla~\ref{tab:e1} la amplía con las demás métricas.

\begin{table}[t]
\centering
\caption{Comparación BFS frente a DFS en la instancia clásica ($N=3$, $K=2$).}
\label{tab:principal}
\setlength{\tabcolsep}{8pt}
\begin{tabular}{@{}lrrr@{}}
\toprule
Algoritmo & Nodos expandidos & Tiempo de ejecución (s) & Longitud del camino \\
\midrule
BFS (cola) & 14 & $2.65\times10^{-5}$ & 11 \\
DFS (pila) & 11 & $2.46\times10^{-5}$ & 11 \\
\bottomrule
\end{tabular}
\end{table}

\begin{table}[t]
\centering
\caption{Métricas completas para $N=3$, $K=2$. Exp.: nodos expandidos; Gen.:
nodos generados; $|F|_{\max}$: tamaño máximo de la frontera; Mem.: máximo de
nodos simultáneos y pico de memoria (\texttt{tracemalloc}); $d$: longitud del
camino. Debajo de la línea, variantes sin detección de estados repetidos.}
\label{tab:e1}
\setlength{\tabcolsep}{4pt}
\resizebox{\textwidth}{!}{\begin{tabular}{@{}lrrrrrrrr@{}}
\toprule
Algoritmo & Exp. & Gen. & $|F|_{\max}$ & Mem. (nodos) & Mem. (KiB) & Tiempo (ms) & $d$ & $b^*$ \\
\midrule
BFS & 14 & 14 & 3 & 15 & 3.2 & 0.0265 & 11 & 1.040 \\
DFS & 11 & 15 & 5 & 16 & 2.9 & 0.0246 & 11 & 1.051 \\
A$^*$ & 14 & 14 & 3 & 15 & 3.4 & 0.0342 & 11 & 1.040 \\
\midrule
BFS (árbol) & 10\,444 & 24\,278 & 13\,835 & 13\,835 & 3\,971.5 & 23.3096 & 11 & 2.383 \\
DFS (árbol) & 11 & 15 & 5 & 5 & 5.1 & 0.0260 & 11 & 1.051 \\
\bottomrule
\end{tabular}}
\end{table}

Los dos algoritmos ciegos encuentran un camino de longitud $d=11$, que es el
óptimo conocido. DFS expande menos nodos (11 frente a 14) porque su primera rama
llega directamente a la meta, mientras que BFS completa cada nivel, incluido el
callejón $\st{3}{2}{0}$ y los estados del nivel 10. Las diferencias de tiempo
($26{,}5$ frente a $24{,}6\,\mu$s) son del orden del ruido de medición y no
permiten afirmar que uno sea más rápido. $b^*\approx 1{,}04$ indica que la
búsqueda es casi lineal. El contraste aparece en las variantes de árbol: sin
detectar repetidos, BFS expande 10\,444 nodos, mantiene 13\,835 en la frontera,
usa 3{,}9\,MiB y tarda 23\,ms (unas 880 veces más que BFS de grafo), con un
$b^*$ aparente de 2{,}38; DFS de árbol, con control de ciclos en el camino,
conserva 5 nodos.

\subsubsection{Sensibilidad de DFS al orden de las acciones.}
El camino que devuelve DFS depende del orden en que se prueban las acciones.
Se ejecutó DFS con todas las permutaciones de $A$
(Tabla~\ref{tab:ordenes}). En la instancia clásica las 120 permutaciones dan
$d=11$; con $N=5$, $K=3$ ($|A|=9$, $9!=362\,880$ órdenes) solo el $16{,}5$\,\%
encuentra el óptimo.

\begin{table}[t]
\centering
\caption{Longitud del camino de DFS sobre todas las permutaciones de las acciones.}
\label{tab:ordenes}
\setlength{\tabcolsep}{6pt}
\begin{tabular}{@{}rrrcrrrr@{}}
\toprule
 & & & & \multicolumn{3}{c}{Órdenes con DFS de longitud} & \\
\cmidrule(lr){5-7}
$N$ & $K$ & Permutaciones & $d^{\,*}$ (BFS) & $d=11$ & $d=13$ & $d=15$ & \% óptimo \\
\midrule
3 & 2 & 120 & 11 & 120 & 0 & 0 & 100{,}0 \\
5 & 3 & 362\,880 & 11 & 59\,832 & 43\,632 & 259\,416 & 16{,}5 \\
\bottomrule
\end{tabular}
\end{table}

\subsubsection{Escalabilidad.}
La Tabla~\ref{tab:e3} y la Fig.~\ref{fig:escala} varían $N$ y $K$ (con $K=2$ no
existe solución para $N\ge 4$, lo que ambos algoritmos detectan al agotar el
espacio). $|\Sr|$ crece linealmente ($6N-2$ para estas instancias) y los nodos
expandidos por los tres algoritmos lo siguen. A medida que el problema crece, DFS
expande algo menos que BFS, pero su camino se aleja del óptimo hasta duplicarlo
($d=1\,015$ frente a $499$ con $N=500$). La frontera de BFS se mantiene en
$|F|_{\max}\le 14$, mientras que la de DFS crece con $N$ hasta 2\,520 nodos.
A$^*$ coincide siempre con el óptimo de BFS, pero ahorra pocas expansiones
(2\,014 frente a 2\,025 con $N=500$).

\begin{table}[t]
\centering
\caption{Escalabilidad de las búsquedas de grafo. $|\Sr|$: estados alcanzables;
Exp.: nodos expandidos; $d$: longitud del camino; ms: mediana del tiempo.}
\label{tab:e3}
\setlength{\tabcolsep}{3.5pt}
\resizebox{\textwidth}{!}{\begin{tabular}{@{}rrr|rrr|rrr|rrr@{}}
\toprule
 & & & \multicolumn{3}{c|}{BFS} & \multicolumn{3}{c|}{DFS} & \multicolumn{3}{c}{A$^*$} \\
$N$ & $K$ & $|S_r|$ & Exp. & $d$ & ms & Exp. & $d$ & ms & Exp. & $d$ & ms \\
\midrule
3 & 2 & 16 & 14 & 11 & 0.026 & 11 & 11 & 0.024 & 14 & 11 & 0.035 \\
5 & 3 & 28 & 26 & 11 & 0.072 & 16 & 15 & 0.054 & 23 & 11 & 0.079 \\
10 & 4 & 58 & 56 & 17 & 0.223 & 36 & 29 & 0.172 & 48 & 17 & 0.225 \\
25 & 4 & 148 & 118 & 47 & 0.498 & 96 & 59 & 0.459 & 108 & 47 & 0.528 \\
50 & 5 & 298 & 223 & 93 & 1.282 & 193 & 111 & 1.240 & 209 & 93 & 1.356 \\
100 & 5 & 598 & 423 & 193 & 2.499 & 393 & 211 & 2.561 & 409 & 193 & 2.727 \\
200 & 6 & 1198 & 825 & 199 & 6.419 & 795 & 415 & 6.836 & 814 & 199 & 7.046 \\
500 & 6 & 2998 & 2025 & 499 & 17.081 & 1995 & 1015 & 18.840 & 2014 & 499 & 19.201 \\
\bottomrule
\end{tabular}}
\end{table}

\begin{figure}[t]
\centering
\begin{tikzpicture}
\begin{axis}[
  name=izq, width=0.53\textwidth, height=4.6cm, xmode=log, ymode=log,
  xlabel={$N$}, ylabel={$|F|_{\max}$ (nodos)}, title={\small Frontera máxima},
  legend style={font=\tiny, at={(0.03,0.97)}, anchor=north west},
  tick label style={font=\scriptsize}, label style={font=\scriptsize},
  log basis x=10]
  \addplot[mark=o, blue] coordinates {(3,3) (5,5) (10,8) (25,9) (50,12) (100,12) (200,14) (500,14)};
  \addplot[mark=square, red] coordinates {(3,5) (5,20) (10,47) (25,67) (50,173) (100,323) (200,1020) (500,2520)};
  \addplot[mark=triangle, black!60!green] coordinates {(3,3) (5,5) (10,9) (25,9) (50,56) (100,107) (200,248) (500,608)};
  \legend{BFS, DFS, A$^*$}
\end{axis}
\begin{axis}[
  at={(izq.east)}, anchor=west, xshift=1.1cm,
  width=0.53\textwidth, height=4.6cm, xmode=log, ymode=log,
  xlabel={$N$}, ylabel={$d$ (cruces)}, title={\small Longitud del camino},
  legend style={font=\tiny, at={(0.03,0.97)}, anchor=north west},
  tick label style={font=\scriptsize}, label style={font=\scriptsize}]
  \addplot[mark=o, blue] coordinates {(3,11) (5,11) (10,17) (25,47) (50,93) (100,193) (200,199) (500,499)};
  \addplot[mark=square, red] coordinates {(3,11) (5,15) (10,29) (25,59) (50,111) (100,211) (200,415) (500,1015)};
  \legend{BFS (óptimo), DFS}
\end{axis}
\end{tikzpicture}
\caption{Frontera máxima y longitud del camino en función de $N$ (instancias de
la Tabla~\ref{tab:e3}, ejes logarítmicos).}
\label{fig:escala}
\end{figure}

\begin{table}[t]
\centering
\caption{Memoria de BFS y DFS (búsqueda de grafo) por instancia.}
\label{tab:mem}
\begin{tabular}{@{}rr|rr|rr@{}}
\toprule
 & & \multicolumn{2}{c|}{BFS} & \multicolumn{2}{c}{DFS} \\
$N$ & $K$ & $|F|_{\max}$ & Mem. (KiB) & $|F|_{\max}$ & Mem. (KiB) \\
\midrule
3 & 2 & 3 & 3.2 & 5 & 2.9 \\
5 & 3 & 5 & 5.2 & 20 & 4.9 \\
10 & 4 & 8 & 6.8 & 47 & 10.9 \\
25 & 4 & 9 & 16.6 & 67 & 24.5 \\
50 & 5 & 12 & 22.1 & 173 & 45.7 \\
100 & 5 & 12 & 56.4 & 323 & 116.3 \\
200 & 6 & 14 & 89.4 & 1020 & 265.4 \\
500 & 6 & 14 & 350.3 & 2520 & 983.7 \\
\bottomrule
\end{tabular}
\end{table}

% ===========================================================================
\section{Discusión y conclusión crítica}\label{sec:discusion}

\subsubsection{¿Por qué difieren los resultados asintóticos?}
La teoría \cite{russell2020} da para BFS tiempo y espacio $O(b^{d})$
($O(b^{d+1})$ con el test objetivo al expandir, como aquí), y para DFS tiempo
$O(b^{m})$ y espacio $O(bm)$, donde $b$ es el factor de ramificación, $d$ la
profundidad de la solución más superficial y $m$ la profundidad máxima. Estas
cotas suponen búsqueda en árbol sobre un espacio grande. En este problema hay
dos hechos que cambian el comportamiento observado:
\begin{enumerate}
\item \emph{El espacio es pequeño y está fuertemente podado.} Con detección de
repetidos, ambos algoritmos quedan acotados por $O(|\Sr|+|E|)$ en tiempo y
$O(|\Sr|)$ en memoria \cite{cormen2022}, con $|\Sr|=16$ para $N=3$ y lineal en
$N$ en general. Por eso los dos expanden casi todo el espacio, y sus tiempos son comparables y crecen de forma polinómica moderada con $N$
(Tabla~\ref{tab:e3}), lejos del crecimiento exponencial de la cota de árbol. El
factor $b^*\approx 1{,}04$ significa que los nodos generados ($14$--$15$) apenas
superan la profundidad de la solución ($d=11$): la búsqueda es casi lineal.
\item \emph{La estructura del grafo es casi una cadena.} La Fig.~\ref{fig:grafo}
muestra que solo hay bifurcaciones al principio y al final. La profundidad máxima
$m$ coincide prácticamente con $d$, así que la ventaja teórica de DFS en memoria
($bm$ frente a $b^d$) y su riesgo de caminos largos ($m\gg d$) no tienen espacio
para manifestarse.
\end{enumerate}
Cuando estas condiciones desaparecen, las diferencias asintóticas vuelven: sin
detección de repetidos el espacio efectivo pasa a ser el árbol de búsqueda, y BFS
expande 10\,444 nodos con $b^*=2{,}38$, mientras que al ampliar $N$ y $K$ el
grafo deja de ser una cadena y DFS empieza a devolver caminos cada vez más largos que el óptimo.

\subsubsection{¿DFS garantizó el camino más corto?}
No. En la instancia clásica DFS devolvió $d=11$ con cualquier orden de las
acciones, pero porque \emph{todos los caminos simples} de $s_0$ a la meta miden
exactamente 11 (verificado por enumeración exhaustiva). Fue una consecuencia de
la estructura del problema, no una propiedad del algoritmo. DFS no es óptimo:
devuelve el primer camino que encuentra su rama más profunda. Con $N=5$, $K=3$,
el $83{,}5$\,\% de los órdenes produce caminos de 13 o 15 cruces, y con $N=500$
el camino de DFS es 2{,}03 veces el óptimo. BFS, en cambio, es óptimo porque los
costos son unitarios: explora por niveles, así que la primera meta extraída está a
profundidad mínima.

\subsubsection{¿Por qué BFS consumió más memoria?}
BFS debe mantener toda la frontera de un nivel antes de pasar al siguiente, y el
tamaño de un nivel crece como $b^{d}$; además, para reconstruir caminos, cada nodo
guarda su padre. DFS solo necesita la rama actual y los hermanos pendientes,
$O(bm)$. Los datos lo confirman en el régimen donde la teoría aplica: en búsqueda de
árbol BFS llega a 13\,835 nodos en la frontera ($\approx$3{,}9\,MiB) frente a 5
nodos de DFS, una razón de 2\,767:1. Es importante, sin embargo, ser precisos con
la búsqueda de grafo de este problema: allí la memoria la domina el conjunto de
explorados, de tamaño $|\Sr|$ en ambos casos, y la diferencia es mínima
(15 frente a 16 nodos; 3{,}2 frente a 2{,}9\,KiB). Incluso, al crecer $N$, la
frontera de DFS supera a la de BFS (Tabla~\ref{tab:mem}), porque la pila conserva
hermanos y duplicados que aún no ha descartado, mientras que el ancho de cada
nivel de BFS está acotado por la propia restricción de seguridad. La afirmación
\emph{``BFS consume más memoria''} es correcta como cota del peor caso
($O(b^d)$ frente a $O(bm)$), pero su magnitud real depende del factor de
ramificación efectivo y de si se eliminan los estados repetidos.

\subsubsection{Conclusión.}
La formulación $(m,c,b)$ con poda en la función de sucesores reduce el problema a
16 estados alcanzables, de modo que BFS y DFS lo resuelven en microsegundos.
Para este tipo de problema, pequeño, con costos unitarios y donde se requiere el
camino mínimo, BFS con detección de repetidos es la elección correcta: garantiza
optimalidad y su costo en memoria queda acotado por $|\Sr|$. DFS solo es
preferible si la memoria es el recurso crítico y no se exige optimalidad. En ese
caso, la profundización iterativa \cite{korf1985} combina la memoria $O(bd)$
de DFS con la optimalidad de BFS. La heurística relajada de A$^*$ ahorra pocas
expansiones porque ignora justo la restricción que hace difícil el problema. Una
heurística más informada debería incorporar la condición de seguridad.

\medskip
\noindent\textbf{Código fuente:}\\ \url{\repo}

\begin{thebibliography}{1}
\providecommand{\url}[1]{\texttt{#1}}
\providecommand{\urlprefix}{URL }
\providecommand{\doi}[1]{https://doi.org/#1}

\bibitem{amarel1968}
Amarel, S.: On representations of problems of reasoning about actions. In:
  Michie, D. (ed.) Machine Intelligence 3, pp. 131--171. Edinburgh University
  Press (1968)

\bibitem{cormen2022}
Cormen, T.H., Leiserson, C.E., Rivest, R.L., Stein, C.: Introduction to
  Algorithms. MIT Press, Cambridge, MA, 4 edn. (2022)

\bibitem{hart1968}
Hart, P.E., Nilsson, N.J., Raphael, B.: A formal basis for the heuristic
  determination of minimum cost paths. IEEE Transactions on Systems Science and
  Cybernetics  \textbf{4}(2),  100--107 (1968)

\bibitem{korf1985}
Korf, R.E.: Depth-first iterative-deepening: An optimal admissible tree search.
  Artificial Intelligence  \textbf{27}(1),  97--109 (1985)

\bibitem{pearl1984}
Pearl, J.: Heuristics: Intelligent Search Strategies for Computer Problem
  Solving. Addison-Wesley, Reading, MA (1984)

\bibitem{russell2020}
Russell, S., Norvig, P.: Artificial Intelligence: A Modern Approach. Pearson,
  Hoboken, NJ, 4 edn. (2020)

\end{thebibliography}

\end{document}
